\chapter{Introduction}
\label{ch:introduction}

\section{Background and Economic Motivation}
Why do stock prices move differently from one another, and can we predict these movements? Economists have tried to answer this question for over 50 years \citep{fama1970efficient, fama1993common}. Classical financial models assume that stock returns depend linearly on a few static risk factors. These include the Capital Asset Pricing Model (CAPM) by \citet{sharpe1964capital} and \citet{lintner1965valuation}, the Arbitrage Pricing Theory by Ross \citeyearpar{ross1976arbitrage}, and the 5-factor model by Fama and French \citeyearpar{fama2015five}.

Real stock markets are much more complex. Stock prices show sudden regime shifts, fat-tailed risks, and non-linear patterns. For example, valuation metrics like Price-to-Earnings ratios behave differently during market crashes than during bull runs \citep{gu2020empirical, kozak2020shrinking}. Standard linear regressions struggle when tracking dozens of stock characteristics at once, leading to overfitting or missed signals.

Machine learning offers a flexible way to capture these complex interactions \citep{gu2020empirical, chen2023deep, kozak2020shrinking}. Models like decision trees \citep{breiman2001random, chen2016xgboost}, neural networks \citep{bai2018empirical}, and Transformers \citep{vaswani2017attention} can learn patterns directly from financial data.

However, stock data is very noisy. Naive machine learning models often fail because they accidentally use future data, overfit historical noise, or ignore trading fees \citep{harvey2016and, novy2016taxonomy, lopez2018advances}. In this thesis, I test whether machine learning models can reliably predict stock returns when tested under strict, real-world trading rules across S\&P 500 stocks from 2015 to 2024.

\section{Research Questions and Core Hypotheses}
I organize my research around four central questions:

\noindent \textbf{Research Question 1 (Predictability)}: Do deep learning models (LSTM, TFDMGA) predict stock returns better than standard linear regressions (OLS, LASSO) and decision trees (Random Forest, XGBoost)?
\begin{quote}
\textit{Hypothesis 1}: Deep sequence models outperform linear models because they track historical time trends and non-linear feature interactions directly.
\end{quote}

\vspace{0.2cm}
\noindent \textbf{Research Question 2 (Multimodal Data Integration)}: Does combining price technicals, financial ratios, macro data, and news sentiment improve forecast stability?
\begin{quote}
\textit{Hypothesis 2}: Combining multiple data types allows attention mechanisms to reweight features as market conditions change, improving signal stability.
\end{quote}

\vspace{0.2cm}
\noindent \textbf{Research Question 3 (Trading Profits \& Costs)}: Can these model predictions build profitable trading strategies after deducting 10 bps trading fees and 1-day execution delays?
\begin{quote}
\textit{Hypothesis 3}: Even though daily rebalancing incurs fee drag, pairing deep learning models with a 2:1 Take-Profit/Stop-Loss rule protects capital and compounds returns.
\end{quote}

\vspace{0.2cm}
\noindent \textbf{Research Question 4 (Market Alpha vs. Factor Exposure)}: Are the extra trading returns coming from unpriced market alpha ($\alpha > 0$) or from systematic factor risks ($\alpha = 0$)?
\begin{quote}
\textit{Hypothesis 4}: The strategy returns are fully explained by standard Fama-French risk factors, showing that machine learning models act as dynamic factor managers rather than finding free money.
\end{quote}

\section{Summary of Key Empirical Findings}
Using a 5-year rolling test timeline from 2020 to 2024, my key empirical findings are:

\begin{enumerate}
    \item \textbf{Top Forecasting Model}: My custom \textbf{TFDMGA} neural network achieves the highest predictive accuracy across all models. It records an out-of-sample daily rank correlation of \textbf{$IC = +0.0348$}, an Information Ratio of \textbf{$ICIR = 3.12$}, a ROC AUC of \textbf{$0.6120$}, and a directional accuracy of \textbf{$56.84\%$}.
    
    \item \textbf{Statistical Comparison}: A formal \citet{diebold1995comparing} test confirms that TFDMGA significantly outperforms a standard PyTorch LSTM model ($DM = 2.41, p = 0.016$).

    \item \textbf{Account Compounding}: Under 10 bps trading fees and 1-day execution delays, an initial \$1,000 USD investment in 2020 grows to \textbf{\$3,120.99 USD} with a baseline LSTM strategy, and \textbf{\$4,368.50 USD} with a 2:1 Stop-Loss rule. Under TFDMGA with the Stop-Loss rule, capital grows to \textbf{\$6,482.10 USD} by 2024.

    \item \textbf{Factor Spanning}: Regressions against the Fama-French 5-factor model yield a net strategy alpha of \textbf{$\hat{\alpha} = -0.18\%$} per year ($p = 0.976$). Because net alpha is zero, 100\% of the returns come from dynamic risk exposure to operating profitability ($RMW$ beta $= +0.512, t = +9.12$), matching market efficiency \citep{fama1970efficient}.
\end{enumerate}

\section{Structure of the Thesis Manuscript}
The remainder of this thesis is structured across four consolidated chapters:
\begin{itemize}
    \item \textbf{Chapter 2 (Literature Review)} surveys academic literature in empirical asset pricing, factor models, high dimensional regularization, and financial machine learning.
    \item \textbf{Chapter 3 (Data \& Methodology)} details the Bloomberg point-in-time data pipeline, constituent panel construction, 59-variable feature taxonomy, baseline model formulations, and the custom TFDMGA network architecture.
    \item \textbf{Chapter 4 (Empirical Results \& Backtesting)} presents daily Fama-MacBeth baselines, LASSO shrinkage derivations, SHAP interpretability analysis, out-of-sample predictive metrics, ablation results, trading backtests under transaction fees, and Fama-French spanning regressions.
    \item \textbf{Chapter 5 (Conclusion \& Future Research Extensions)} summarizes the main findings, discusses theoretical and practical implications, addresses study limitations, and outlines future research extensions.
\end{itemize}
